\section{Method}
\label{sec:method}

\subsection{Dataset}
\label{sec:dataset}

We evaluate on \textbf{ConTRoL} \citep{liu-etal-2021-natural}, a dataset for
document-level natural language inference constructed from expert-designed
reading-comprehension material. Each ConTRoL instance pairs a long premise
(averaging approximately 500 words) with a short hypothesis, and requires
identifying whether the hypothesis is \textbf{Entailed}, \textbf{Contradicted},
or \textbf{Neutral} with respect to the premise. The dataset comprises 8,325
premise-hypothesis pairs. Premises are long, multi-sentence passages rather
than single sentences, so the evidence needed to decide a label is often
spread across several sentences rather than contained in one. This makes
ConTRoL more demanding than sentence-level NLI benchmarks.

\subsection{Sampling}
\label{sec:sampling}

All experiments were run against fixed samples drawn from the ConTRoL train
split, using a fixed random seed (seed~$=$~42).\footnote{A na\"ive
lexicographic sort over string-formatted instance IDs (e.g.\ \texttt{id\_1},
\texttt{id\_10}, \texttt{id\_100}, \ldots, \texttt{id\_11}) does not
correspond to a numeric or random ordering and yields a non-representative
sample.} The sampled instance IDs were persisted once and reused identically
across every model, ensuring that all models and methods are evaluated on the
exact same data.

Due to the substantially higher computational cost of the premise-decomposition
methods (Section~\ref{sec:methods}), we adopt two sample sizes: $N=1{,}500$
for the six methods that issue a bounded number of model calls per instance,
and $N=150$ for the three premise-question configurations, whose per-instance
cost scales with the number of decomposed sub-questions.

\subsection{Seeding Strategies}
\label{sec:seeding}

Two of our pipelines (Section~\ref{sec:methods}) generate probe questions from
an anchor set of extracted spans, rather than directly from the raw text. We
compare two seeding strategies for producing this anchor set.

\paragraph{Part-of-Speech (POS) seeding.} We use NLTK's POS tagger
\citep{loper2002nltk} to extract noun phrases, main verbs, and numeric spans
from the source text via shallow chunking. This method requires no model
inference and is fully deterministic.

\paragraph{Semantic Role Labeling (SRL) seeding.} Following the QA-SRL
paradigm \citep{fitzgerald2018large,pyatkin2021asking}, we prompt the
evaluated LLM to extract PropBank-style predicate-argument structures ---
agent (ARG0), theme (ARG1), recipient (ARG2), and modifier roles for time,
location, manner, and cause --- for each predicate in the source text. Each
role-argument pair is converted into an anchor of the form
\texttt{\{role\}: \{span\} [predicate: \{verb\}]}. This strategy is more
computationally expensive than POS seeding, requiring one additional model
call, but yields anchors with explicit predicate-argument structure rather
than flat lexical spans.

\subsection{Methods}
\label{sec:methods}

We evaluate nine method configurations, organized into three groups.

\paragraph{Direct baselines.} \texttt{zero-shot} issues the premise and
hypothesis to the model in a single call with no intermediate reasoning step.
\texttt{few-shot-CoT} uses the same single-call structure, prepending three
fixed demonstrations --- one worked example per label, held constant across
all evaluations --- following \citet{wei2022chain}.

\paragraph{Retrieval control.} \texttt{retrieve-then-classify} isolates the
contribution of evidence localization from that of question generation. A
locator model identifies the premise sentences most relevant to the
hypothesis, using the \textbf{hypothesis itself} as the retrieval query; the
classifier is then given these sentences verbatim, without an intermediate
answer-extraction step. This condition allows us to attribute any additional
gain from the question-based methods below specifically to the generated
questions, rather than to evidence retrieval alone.

\paragraph{Question-based pipelines.} The remaining six configurations share a
common three-stage architecture: (i) question generation, (ii) evidence
localization and answer extraction, and (iii) classification of the
hypothesis given the resulting (question, answer) pairs. Stages (ii) and
(iii) are implemented identically across all six configurations (detailed
below), so that any difference in downstream accuracy can be attributed to
\emph{which} questions were generated and \emph{which text they were
generated from}, not to how they are answered or judged. The three
question-based methods differ precisely along this second axis --- questions
generated from the hypothesis, from the premise, or from both --- letting us
compare which side of the premise-hypothesis pair is more productive to
interrogate.

\subsubsection{Answering and Classification}

\paragraph{Locate and answer.} A locator model reads the premise and one
question and returns up to five relevant sentence indices (not necessarily
adjacent, allowing multi-hop evidence). An extractor model then answers the
question using only those sentences, returning a not-answerable marker
instead of guessing if they do not address it; such questions are dropped.

\paragraph{Classify.} The surviving (question, answer) pairs and the
hypothesis are given to the classifier in one call, which returns a single
Entailment / Contradiction / Neutral label. Priority is Contradiction over
Entailment over Neutral, and one clearly decisive piece of evidence is
treated as sufficient on its own, without being diluted by other, merely
tangential evidence in the same call.

Stage (i) --- question generation --- is the only stage that differs between
the three methods below.

\begin{itemize}
    \item \textbf{h-question} generates 2--4 probe questions directly from
    the hypothesis, seeded via either POS or SRL extraction
    (Section~\ref{sec:seeding}) applied to the hypothesis text. This pipeline
    is hypothesis-aware but premise-blind at generation time.

    \item \textbf{bridge-question} generates 2--4 sub-questions from the
    premise and hypothesis jointly, with at least one question explicitly
    required to test the hypothesis's core claim against premise evidence,
    following the bridging/comparison framing of multi-hop QA
    \citep{yang-etal-2018-hotpotqa,min-etal-2019-multi}. This is the only configuration
    whose question-generation stage observes both texts simultaneously.

    \item \textbf{p-question (decomposition)} applies \emph{atomic
    decomposition} to the premise: the model is prompted to break the
    premise into the smallest standalone factual claims it contains ---
    each covering exactly one date, name, number, or event --- following
    the atomic-fact paradigm of FActScore \citep{min-etal-2023-factscore}. A claim
    is ``atomic'' in the sense that it cannot be split further without
    losing a checkable unit of information; a premise stating that ``the
    government paid \$1~lakh to families of railway-accident victims''
    atomizes into separate facts for the amount, the recipients, and the
    triggering event. In a second step, the model additionally identifies
    \emph{relations} between these facts --- comparisons, causal links, and
    temporal orderings --- following relation-level verification in the
    style of \citet{goyal2020evaluating,goyal2021annotating}. Each generated
    question is tagged as targeting either a single fact or a relation
    between facts. This two-step decomposition guarantees that both the
    premise's individual claims and the connections between them are
    covered by construction, rather than left to the model's unconstrained
    judgment of what is salient. This mode is premise-aware but
    hypothesis-blind at generation time.

    \item \textbf{p-question (seeded)} generates one question per anchor
    extracted from the premise via the seeding strategies of
    Section~\ref{sec:seeding}, evaluated separately for POS and SRL
    seeding.
\end{itemize}

Stage-1 question generation additionally receives three worked examples
in-context for all six question-based configurations, following the same
few-shot design as \texttt{few-shot-CoT}.

\subsection{Models}
\label{sec:models}

We evaluate four models spanning both self-hosted open-weight and API-hosted
settings: \textbf{Qwen2.5-32B} and \textbf{gpt-oss-20B}, run locally via
Ollama-backed infrastructure, and \textbf{GPT-4o-mini} and
\textbf{DeepSeek-V4-Flash}, accessed via their respective provider APIs. This
selection allows comparison across model scale and serving infrastructure,
while controlling the evaluation pipeline, prompts, and sampled data
identically across all four.
